%%
%% Author: shoupu_wan
%% 2019-05-22
%%

\section{Sentinels in quicksort}\label{sec:quick-sort-sentinel}
%TODO Simplify the presentation of sentinels to one to two sentences

The cascading conditional in the loop body of ~\autoref{lst:singleLoopHoarePart}
is still far from straightforward.
A further simplification on that may achievable through the use of
sentinels~\cite{NumericalRecipes:1992}.

%It was eye-opening to see how the authors gracefully avoided the array boundary checks
%with the use of sentinels.
%I psychologically accepted the pitch made for sentinels while I was with the book.
%Closing the book, however, I am still dubious about its plausibility until I myself have a
%success implementing it.
%In this section, I pick this work to make a pitch to the reader about sentinels.


For stark effect of the sentinels, we will first fix the traditional implementation of `Hoare's
partition scheme':
~\autoref{lst:QuickSortOutline} gives the outline of the program and
~\autoref{lst:nestedLoopHoareMethod} gives a fully implemented partition function.
As mentioned earlier, one can get involved in the thick of implementing
`Hoare's quicksort algorithm'.
Not only the problem itself is tricky but also the way we implement it.
By picking Hoare's quicksort scheme over its alternative (such as Lomuto's),
we insist on the robust $O(N log N)$ expected runtime.


We attempt to get rid of the boundary checking using sentinels.
These boundary checking are there to ensure that the pointers `\code{s}' and `\code{e}' do not slip
off the ends of the array.
In large arrays, these boundary check operations may get in the way of performance of the algorithm.
But the main reason of doing so is to remove clutter in the code.
%For another, we are consciously attacking the problem with an precautions approach against
%performance degradation so there is a price we have to pay.

The main idea is to make boundary checking redundant by effecting some artifacts that catch the
pointers before the `out of boundary' error happens.
This is exactly what sentinel is good at!
By deploying sought-after values, the `sentinels', at the ends of the array before the control
hits the loop, the following will happen:
\begin{enumerate}
    \item the pointers, `\code{s}' and `\code{p}', would have to stop when they hit the ends of
    array;
    \item the subsequent logic will guarantee the termination of the outer loop and
    thus guarantee the inner loop would not be executed again.
\end{enumerate}


What are the sought-after values by `\code{s}' and `\code{e}'?
The out-of-place values!
Particularly, \code{s} is looking for values that are `$\geq$ the pivot';
\code{e} is looking for values `$\leq$ the pivot'.
So then instead of randomly picking one value for pivot, we now randomly pick $3$ values, the median
of which be elected as the pivot as usual, the minimum be deployed to the left end, and the maximum
to the right end.
We group these operations into a function called `\code{init_swap}':
\lstinputlisting[]{qsort_init_swap.c}


Now back the function `\code{part}'.
As we mentioned before, the deployment of sentinels in function `\code{init_swap}' makes
the conditional expressions `\code{s < e}' and `\code{pivot < e}' semantically redundant.
They can now be safely removed.
Below is the function `\code{part}' after all the refactoring is done.
\lstinputlisting[]{qsort_part_sentinel.c}
Comparing with the implementation ~\autoref{lst:nestedLoopHoareMethod}, this code
cleans out clutter in the loop conditions in the inner loops.
We `float' the control all the way through the termination of the program on a `touchless' rail made
by sentinels, without ever needing to check boundary condition.
Of course, the onus is shifted to the initialization before the loop.
Comparing with the inner loops, that section is non-critical.
Not only that the code is made less cluttered, but also the number of such checks
is reduced from $O(N)$ to $O(1)$.
More importantly, by not checking boundary conditions at the most critical section,
we avoid bugs that would otherwise cost us many hours of debugging time.


So by use of sentinels on the traditional Hoare's quicksort algorithm,
we shifted the complexity among the nested loops to a non-critical part of the program,
effectively reduced its coding complexity.
Remember that we came up with a new, single-loop implementation for the Hoare's quicksort algorithm
in ~\autoref{ch:loop-thinning}, as presented in ~\autoref{lst:singleLoopHoarePart}.
Can we apply the idea of sentinel to make it even better?
Yes, it may be applied equally well.
Below is the implementation of the function `\code{part}' with a single loop and simplified
boundary condition with sentinel values.
\lstinputlisting[]{quicksort_partition_single_loop_sentinel.c}
This brings implementations of Hoare's quicksort to the next level of simplicity.
